% GENERATED FILE. Edit canonical lessons, then run npm run generate:books.
% canonical-source-hash: fnv1a64:b162f05eb81a41ed
% canonical-lessons: KA-C251-gadi, KA-C251-kaluve, KA-C251-suranga, KA-C251-ilijaru, KA-C251-hittilu

\chapter{Frontier, Canal, Tunnel, Downward slope, Backyard}
\label{ch:ka-frontier-canal-tunnel-downward-slope-backyard}
\hlchaptermodality{251}

\begin{chapteropening}
\textbf{By the end of this chapter:} \emph{I can say a frontier, a canal, a tunnel, a downward slope and a backyard.}

Complete the last lesson of chapter 251: I can say a frontier, a canal, a tunnel, a downward slope and a backyard.
\end{chapteropening}

\section[\texorpdfstring{\kn{ಗಡಿ} (gaḍi)}{ಗಡಿ (gaḍi)}]{\kn{ಗಡಿ} (gaḍi) — a frontier, a boundary}

\label{lesson:KA-C251-gadi}



\begin{quote}
Before the new one: say the Kannada for a post office, then the Kannada for a valley.
\end{quote}

\subsection*{You'll want to know: \kn{ಗಡಿ}}
\textbf{\kn{ಗಡಿ}} — \emph{gaḍi} — \textquotedblleft{}a frontier, a boundary\textquotedblright{}.

\begin{grammarlens}[title={What you've built}]
One more word for this chapter.
\end{grammarlens}

\subsection*{Your turn}
\begin{itemize}
\raggedright
  \item \emph{Say it:} \emph{gaḍi}
  \item \emph{Say it:} \emph{gaḍi}, once more
  \item \emph{Say it:} say \emph{gaḍi}
  \item \emph{Recall:} say the Kannada for a post office, then the Kannada for a valley, then say \emph{gaḍi} again
\end{itemize}

\begin{tcolorbox}[breakable,colback=teal!4,colframe=teal!35!black,title={Before you move on}]
What does \kn{ಗಡಿ} mean? (\textquotedblleft{}A frontier, a boundary\textquotedblright{}.) Say it once more.
\end{tcolorbox}
\section[\texorpdfstring{\kn{ಕಾಲುವೆ} (kāluve)}{ಕಾಲುವೆ (kāluve)}]{\kn{ಕಾಲುವೆ} (kāluve) — a canal, a channel}

\label{lesson:KA-C251-kaluve}



\begin{quote}
Before the new one: say the Kannada for a valley, then the Kannada for a frontier.
\end{quote}

\subsection*{You'll want to know: \kn{ಕಾಲುವೆ}}
\textbf{\kn{ಕಾಲುವೆ}} — \emph{kāluve} — \textquotedblleft{}a canal, a channel\textquotedblright{}.

\begin{grammarlens}[title={What you've built}]
One more word for this chapter.
\end{grammarlens}

\subsection*{Your turn}
\begin{itemize}
\raggedright
  \item \emph{Say it:} \emph{kāluve}
  \item \emph{Say it:} \emph{kāluve}, once more
  \item \emph{Say it:} say \emph{kāluve}
  \item \emph{Recall:} say the Kannada for a valley, then the Kannada for a frontier, then say \emph{kāluve} again
\end{itemize}

\begin{tcolorbox}[breakable,colback=teal!4,colframe=teal!35!black,title={Before you move on}]
What does \kn{ಕಾಲುವೆ} mean? (\textquotedblleft{}A canal, a channel\textquotedblright{}.) Say it once more.
\end{tcolorbox}
\section[\texorpdfstring{\kn{ಸುರಂಗ} (suraṅga)}{ಸುರಂಗ (suraṅga)}]{\kn{ಸುರಂಗ} (suraṅga) — a tunnel}

\label{lesson:KA-C251-suranga}



\begin{quote}
Before the new one: say the Kannada for a frontier, then the Kannada for a canal.
\end{quote}

\subsection*{You'll want to know: \kn{ಸುರಂಗ}}
\textbf{\kn{ಸುರಂಗ}} — \emph{suraṅga} — \textquotedblleft{}a tunnel\textquotedblright{}.

\begin{grammarlens}[title={What you've built}]
One more word for this chapter.
\end{grammarlens}

\subsection*{Your turn}
\begin{itemize}
\raggedright
  \item \emph{Say it:} \emph{suraṅga}
  \item \emph{Say it:} \emph{suraṅga}, once more
  \item \emph{Say it:} say \emph{suraṅga}
  \item \emph{Recall:} say the Kannada for a frontier, then the Kannada for a canal, then say \emph{suraṅga} again
\end{itemize}

\begin{tcolorbox}[breakable,colback=teal!4,colframe=teal!35!black,title={Before you move on}]
What does \kn{ಸುರಂಗ} mean? (\textquotedblleft{}A tunnel\textquotedblright{}.) Say it once more.
\end{tcolorbox}
\section[\texorpdfstring{\kn{ಇಳಿಜಾರು} (iḷijāru)}{ಇಳಿಜಾರು (iḷijāru)}]{\kn{ಇಳಿಜಾರು} (iḷijāru) — a downward slope}

\label{lesson:KA-C251-ilijaru}



\begin{quote}
Before the new one: say the Kannada for a canal, then the Kannada for a tunnel.
\end{quote}

\subsection*{You'll want to know: \kn{ಇಳಿಜಾರು}}
\textbf{\kn{ಇಳಿಜಾರು}} — \emph{iḷijāru} — \textquotedblleft{}a downward slope\textquotedblright{}.

\begin{grammarlens}[title={What you've built}]
One more word for this chapter.
\end{grammarlens}

\subsection*{Your turn}
\begin{itemize}
\raggedright
  \item \emph{Say it:} \emph{iḷijāru}
  \item \emph{Say it:} \emph{iḷijāru}, once more
  \item \emph{Say it:} say \emph{iḷijāru}
  \item \emph{Recall:} say the Kannada for a canal, then the Kannada for a tunnel, then say \emph{iḷijāru} again
\end{itemize}

\begin{tcolorbox}[breakable,colback=teal!4,colframe=teal!35!black,title={Before you move on}]
What does \kn{ಇಳಿಜಾರು} mean? (\textquotedblleft{}A downward slope\textquotedblright{}.) Say it once more.
\end{tcolorbox}
\section[\texorpdfstring{\kn{ಹಿತ್ತಿಲು} (hittilu)}{ಹಿತ್ತಿಲು (hittilu)}]{\kn{ಹಿತ್ತಿಲು} (hittilu) — a backyard}

\label{lesson:KA-C251-hittilu}



\begin{quote}
Before the new one: say the Kannada for a tunnel, then the Kannada for a downward slope.
\end{quote}

\subsection*{You'll want to know: \kn{ಹಿತ್ತಿಲು}}
\textbf{\kn{ಹಿತ್ತಿಲು}} — \emph{hittilu} — \textquotedblleft{}a backyard\textquotedblright{}.

\begin{grammarlens}[title={What you've built}]
One more word for this chapter.
\end{grammarlens}

\subsection*{Your turn}
\begin{itemize}
\raggedright
  \item \emph{Say it:} \emph{hittilu}
  \item \emph{Say it:} \emph{hittilu}, once more
  \item \emph{Say it:} say \emph{hittilu}
  \item \emph{Recall:} say the Kannada for a tunnel, then the Kannada for a downward slope, then say \emph{hittilu} again
\end{itemize}

\begin{tcolorbox}[breakable,colback=teal!4,colframe=teal!35!black,title={Before you move on}]
What does \kn{ಹಿತ್ತಿಲು} mean? (\textquotedblleft{}A backyard\textquotedblright{}.) Say it once more.
\end{tcolorbox}
